%% ===================================================================================
%% Capítulo 1 -- Introdução
%% Documento de Arquitetura e Requisitos de Software do Vitrine. Origem do produto:
%% Relatório de IC (doc/ref). Dono: Arquiteto/Redator.
%% ===================================================================================
\chapter{Introdução}\label{cap:intro}

\section{Propósito e Público do Documento}\label{sec:proposito}

Este documento é a especificação oficial do Vitrine: reúne os requisitos, a arquitetura, o modelo de dados, os contratos de API, o plano de implementação e o plano de verificação do sistema. É a \textbf{Fonte da Verdade} do projeto: o código a implementa e o Arquiteto confere a aderência a cada integração (Seção~\ref{sec:processo-ddd}). Ele é mantido junto ao código, em \LaTeX{}, e evolui com o produto.

O público do documento é o Engenheiro Chefe, que decide o rumo do produto; os agentes e desenvolvedores dos domínios Dados, Core e UI, que implementam; o Implantador, que instala e opera o sistema (Apêndice~\ref{ap:implantacao}); e quem quiser estender ou reimplementar o Vitrine.

\section{Visão Geral do Produto}\label{sec:visao-geral-produto}

A disseminação do conhecimento acadêmico encontrou na web o seu principal canal de expansão \cite{ginige2001web}. Docentes e pesquisadores, especialmente os de áreas não tecnológicas, dependem cada vez mais de um espaço próprio para publicar artigos, notas de aula e materiais didáticos, com identidade reconhecível e endereço estável. Entretanto, a criação desse espaço ainda esbarra em um dilema: de um lado, a dependência de plataformas proprietárias e rígidas, que limitam o controle sobre a identidade visual, a formatação de elementos técnicos --- fórmulas matemáticas e blocos de código --- e a perenidade do conteúdo; de outro, a complexidade de instalar, configurar e manter sistemas de gerenciamento de conteúdo (\en{Content Management Systems}, CMS) tradicionais \cite{mikkonen2010mashware}.

O Vitrine é a resposta a esse problema: uma plataforma de gerenciamento de conteúdo de código aberto, sob licença MIT (\adr{ADR-015}), concebida para a criação de blogs e portfólios acadêmicos. Combina um editor Markdown com suporte a fórmulas em \LaTeX{}, uma estrutura hierárquica de documentos e coleções, um estado de publicação por conteúdo e um sistema de personalização visual aplicado em tempo de execução. A primeira versão foi concebida e prototipada em um projeto de Iniciação Científica (IC) entre setembro de 2025 e fevereiro de 2026 \cite{reis2026ic}; esta documentação consolida essa base em uma arquitetura monolítica modular sobre Next.js, SQLite e better-auth.

A premissa de produto é que cada educador deve poder possuir a sua própria ``vitrine'' de conhecimento, de forma independente e personalizada: a implantação deve ser um ato acessível, e a barreira técnica não deve impedir a presença digital qualificada do docente. Essa premissa dialoga com a visão de que sistemas de informação eficazes devem ser facilitadores dos processos humanos e não obstáculos a eles \cite{rainer2019intro}.

\section{Objetivos do Produto}\label{sec:objetivos}

\subsection{Objetivo Geral}

Oferecer uma plataforma CMS de código aberto, modular e de fácil implantação, baseada no \en{framework} Next.js, voltada à criação de blogs e portfólios acadêmicos por professores e pesquisadores, descrita por uma documentação que seja a fonte da verdade que orienta e verifica a sua implementação.

\subsection{Objetivos Específicos}\label{sec:objetivos-especificos}

Os objetivos OE1 a OE5 são herdados do relatório de IC \cite{reis2026ic}; o único ajuste é a adoção da biblioteca better-auth no lugar do Auth.js inicialmente previsto (decisão \adr{ADR-004}).

\begin{enumerate}[label=\textbf{OE\arabic*.}, leftmargin=1.6cm]
  \item Viabilizar a publicação independente por meio de um processo de clonagem e implantação simplificado;
  \item Implementar um editor de texto baseado em Markdown que atenda às necessidades da escrita científica, como realce de sintaxe e fórmulas em \LaTeX{};
  \item Prover um sistema de personalização visual dinâmica que permita alterar a paleta de cores e a identidade do site por meio de interface administrativa, sem novos ciclos de compilação (\en{build});
  \item Adotar uma arquitetura de dados portátil, com o ORM Prisma e o banco de dados SQLite;
  \item Integrar um sistema de autenticação simplificado, por meio do better-auth, para garantir a segurança no gerenciamento de conteúdos.
\end{enumerate}

\section{Justificativa Técnica}\label{sec:justificativa}

O equilíbrio entre modernidade tecnológica e portabilidade orienta as escolhas. O Next.js permite uma renderização híbrida de alto desempenho, essencial para a indexação por motores de busca (\en{Search Engine Optimization}, SEO) \cite{nextjsdocs}. O SQLite, por ser um sistema baseado em arquivo, simplifica a infraestrutura e reduz custos para o usuário final \cite{sqliteabout,sqlitewhentouse}. A camada de acesso a dados com Prisma ORM isola o restante do sistema do dialeto SQL e oferece um modelo declarativo versionável \cite{prismadocs}. As customizações visuais são feitas por variáveis CSS em tempo de execução, com persistência em banco de dados e reflexo imediato por meio de \en{Incremental Static Regeneration} (ISR) \cite{nextjsdocs}.

Embora o Vitrine seja intuitivo para usuários leigos, ele também atende a desenvolvedores, ao oferecer uma base sólida que elimina o trabalho repetitivo de configurar estruturas do zero. O valor do produto pode ser resumido em três dimensões:

\begin{itemize}
  \item \textbf{Social:} contribui para a democratização da publicação acadêmica ao reduzir a dependência de plataformas proprietárias;
  \item \textbf{Prática:} entrega ao docente uma ferramenta pronta para uso, com suporte a fórmulas, código e identidade visual própria;
  \item \textbf{Tecnológica:} demonstra uma arquitetura monolítica enxuta, sem serviços externos obrigatórios, que combina desempenho estático e flexibilidade dinâmica.
\end{itemize}

\section{Escopo e Delimitações}\label{sec:escopo}

O escopo desta versão do Vitrine compreende: autenticação de um administrador; gestão hierárquica de documentos e coleções, com estado de publicação (rascunho ou publicado); edição e renderização de Markdown com fórmulas \LaTeX{}, tabelas e código; navegação por barra superior e \en{sidebar} hierárquica; personalização de identidade, paleta de cores e links sociais do rodapé; e entrega das páginas públicas por ISR com invalidação sob demanda, implantadas por imagem Docker em uma VPS.

Ficam fora do escopo desta versão, embora documentados como possibilidades de evolução (Seção~\ref{sec:trabalhos-futuros}): múltiplos usuários com papéis e permissões distintos; comentários e interação de leitores; busca textual; histórico de versões de conteúdo; envio e gerenciamento de mídia; internacionalização; e edição colaborativa em tempo real.

\section{Histórico do Projeto: da Iniciação Científica à Versão Atual}\label{sec:historico}

O relatório de IC \cite{reis2026ic} registra a concepção do sistema, os requisitos iniciais, os diagramas de casos de uso e de entidade-relacionamento, os protótipos de baixa fidelidade e a decisão de usar Next.js, Prisma e SQLite. Esta documentação reaproveita integralmente esse acervo e o evolui. A Tabela~\ref{tab:ic-vs-atual} resume o que permanece e o que muda.

\begin{table}[htb]
\caption{Comparação entre o relatório de IC e a versão atual}\label{tab:ic-vs-atual}
\small
\begin{tabularx}{\linewidth}{@{}L{2.9cm}YY@{}}
\toprule
\textbf{Aspecto} & \textbf{Relatório de IC (set/2025 a fev/2026)} & \textbf{Versão atual}\\
\midrule
Objetivo & Conceber e prototipar o CMS acadêmico & Consolidar, documentar e verificar o monólito\\
\addlinespace
Autenticação & Auth.js (previsto) & better-auth (\adr{ADR-004})\\
\addlinespace
Modelo de dados & Três entidades: usuários, conteúdos e configurações & Seis modelos: quatro do better-auth (\cod{user}, \cod{session}, \cod{account}, \cod{verification}) mais conteúdos e configurações\\
\addlinespace
Especificação & Relatório em PDF, diagramas exportados & Documentação \LaTeX{} viva e diagramas como código (PlantUML e DBML)\\
\addlinespace
Processo & Incremental, individual & \en{Documentation-Driven Development} com agentes especializados e mural\\
\addlinespace
Interface & Protótipos de baixa fidelidade e uma instância de demonstração & Implementação completa conforme a especificação (Capítulo~\ref{cap:implementacao})\\
\addlinespace
Cache & ISR citado como requisito & ISR com revalidação sob demanda por rota e por \en{layout} (Seção~\ref{sec:arq-isr})\\
\addlinespace
Implantação & Clonagem do repositório e implantação em plataforma como a Vercel (tutoriais previstos) & Imagem Docker em VPS tradicional, com volume persistente para o SQLite (\adr{ADR-011})\\
\addlinespace
Conteúdo & Sem estado de publicação no modelo de dados (o protótipo oferece ``Salvar localmente'' e ``Enviar nova versão'') & Estado de publicação \texttt{DRAFT} ou \texttt{PUBLISHED}, com padrão \texttt{PUBLISHED} (\adr{ADR-012})\\
\bottomrule
\end{tabularx}
\fonte{Elaborado pelos autores, com base em \citeonline{reis2026ic} e no estado do repositório.}
\end{table}

Uma consequência dessa evolução é que alguns termos do relatório de IC --- por exemplo, ``Auth.js'' e a tabela ``\cod{users}'' --- foram substituídos nesta documentação. Sempre que o texto fizer referência ao relatório de IC, a origem será explicitada.

\section{Convenções do Documento}\label{sec:convencoes}

Os elementos do projeto têm identificadores estáveis, citados nos \en{commits}, nos pedidos de integração e nos testes: \texttt{UC} (casos de uso), \texttt{RF} e \texttt{RNF} (requisitos), \texttt{RN} (regras de negócio), \texttt{ADR} (decisões de projeto), \texttt{M} (módulos) e \texttt{CT} (casos de teste); as referências cruzadas são hiperligações. Trechos marcados como \emph{normativos} são contratos a cumprir; os marcados como \emph{esboço} apenas ilustram e não vinculam a implementação. Blocos tracejados amarelos indicam figuras ou trechos de código ainda não produzidos e são listados no Apêndice~\ref{ap:pendencias}.

\section{Organização do Documento}\label{sec:organizacao}

Além desta introdução, o documento está organizado da seguinte forma:

\begin{itemize}
  \item \textbf{Capítulo~\ref{cap:fundamentacao}:} apresenta o contexto tecnológico --- CMS, arquiteturas e estratégias de renderização, tecnologias adotadas, engenharia de requisitos e \en{Documentation-Driven Development};
  \item \textbf{Capítulo~\ref{cap:metodologia}:} descreve a abordagem de desenvolvimento, o planejamento por fases, o processo com agentes, o ambiente e os riscos;
  \item \textbf{Capítulo~\ref{cap:requisitos}:} apresenta o levantamento e a análise de requisitos: atores, casos de uso, requisitos funcionais e não funcionais, regras de negócio e prototipagem;
  \item \textbf{Capítulo~\ref{cap:arquitetura}:} detalha o projeto da arquitetura: visões, decisões, modelo de dados, autenticação, renderização e cache, contrato de API, segurança e implantação;
  \item \textbf{Capítulo~\ref{cap:implementacao}:} descreve a implementação por módulo, com o estado de cada um;
  \item \textbf{Capítulo~\ref{cap:testes}:} define a estratégia e o plano de verificação e validação;
  \item \textbf{Capítulo~\ref{cap:resultados}:} registra a situação da verificação, as limitações conhecidas e o roteiro de evolução.
\end{itemize}

Os apêndices reúnem o dicionário de dados, o esquema Prisma, os fontes dos diagramas, o guia de implantação, os registros de decisão de projeto (ADR), o mapa de pendências da documentação e o glossário.
